\begin{listing}[H]
        \caption{Sample \Feelpp code for solving a Laplace equation on an arbitrary domain.}
\begin{minted}[
        linenos,                % Line numbers
        fontsize=\scriptsize,        % Reduce font size
        bgcolor=bgcolor,        % Slightly gray background
        frame=lines,            % Delimiters around the code
        framesep=2mm,           % Space between code and frame
        rulecolor=\color{gray}, % Color of the frame
        breaklines              % Allow line breaks in long lines
      ]{cpp}
auto Vh = Pch<4>(mesh, markedelements(mesh, expr("<...>")));
auto u = Vh->element(), v = Vh->element(g, "g");
auto l = form1(_test = Vh);
l = integrate(_range = elements(support(Vh)), _expr = f * id(v));
l += integrate(_range = markedfaces(support(Vh), "Robin"), _expr = -r_2 * id(v));
l += integrate(_range = markedfaces(support(Vh), "Neumann"), _expr = -un * id(v));

auto a = form2(_trial = Vh, _test = Vh);
a = integrate(_range = elements(support(Vh)), _expr = inner(k * gradt(u), grad(v)));
a += integrate(_range = markedfaces(support(Vh), "Robin"), _expr = r_1 * idt(u) * id(v));
a += on(_range = markedfaces(support(Vh), "Dirichlet"), _rhs = l, _element = u, _expr = g);
a.solve(_rhs = l, _solution = u);
\end{minted}
\end{listing}